\chapter{Bank Markets Divisions}\label{ch:in:bank-markets-divisions}

Every large bank reports its markets revenue in two lines: fixed income, currencies and commodities, and equities.
The split differs across banks by more than their size does. In 2025 one US bank earned 64\% of its markets revenue in
equities and another 26\%; a French bank earned 61\% in equities while describing its wholesale business as having
``distinctive global leadership in equity derivatives''; a German bank earned none at all, having announced in July 2019
that it would exit equities sales and trading. The mix is the bank's product strength, and it decides which quant jobs
the bank has: a bank strong in equity derivatives employs pricing quants for structured products and volatility, one
strong in rates employs curve and XVA quants. This chapter reads banks' markets divisions as employers through what
they publish.

\section{What banks publish about their markets businesses}

Banks are the most transparent employers in this book. Listed banks file annual reports with segment results
(Book~16, chapter~5), and inside the investment-banking segment most report markets revenue by product line, quarter by
quarter. What they publish, and what they do not, shapes what a candidate can learn.

\begin{definition}[FICC]\label{def:in:bank-markets-divisions:ficc}
\emph{FICC}\index{FICC} (fixed income, currencies and commodities) is the markets business of a bank that deals in
government and corporate bonds, interest-rate and credit derivatives, foreign exchange, commodities and their
financing; its counterpart line is equities: cash shares, equity derivatives and prime services.
\end{definition}

Each bank draws the line its own way. One reports ``fixed income markets'' and ``equity markets''; another ``\omterm{def:in:bank-markets-divisions:ficc}{FICC}'' and
``equities'' inside a segment that also contains advisory; a third reports ``fixed income'' and ``equity'' inside its
institutional securities segment; a European bank reports ``Equity \& Prime Services''; another includes financing
activities in fixed income. Comparisons are therefore of shares within each bank, not of amounts across banks, and
each bank's own line names are kept beside its numbers.

\begin{omfigure}
\begin{tikzpicture}[font=\footnotesize,>=stealth,
  top/.style={draw,rounded corners,fill=omDef!12,minimum width=3.4cm,minimum height=0.6cm,align=center},
  desk/.style={draw,rounded corners,fill=omProp!8,minimum width=4.6cm,minimum height=0.72cm,align=center}]
\node[top] (m) at (4.5,0) {markets division};
\node[top] (f) at (1.5,-1.2) {FICC};
\node[top] (e) at (7.5,-1.2) {equities};
\draw (m.south) -- ++(0,-0.3) -| (f.north); \draw (m.south) -- ++(0,-0.3) -| (e.north);
\foreach \n/\y/\t/\q in {r/-2.2/rates/{curve and rates-option quants},c/-3.15/credit/{credit and XVA quants},
  x/-4.1/foreign exchange/{FX-option and e-FX quants},k/-5.05/commodities/{commodity model quants}}
  {\node[desk] (\n) at (1.5,\y) {\t\\[-1pt]{\scriptsize\itshape \q}};}
\foreach \n/\y/\t/\q in {q/-2.2/cash equities/{execution and index quants},d/-3.15/equity derivatives/{volatility and structuring quants},
  p/-4.1/prime services/{financing and margin quants}}
  {\node[desk] (\n) at (7.5,\y) {\t\\[-1pt]{\scriptsize\itshape \q}};}
\draw (f.south) -- (r.north); \draw (e.south) -- (q.north);
\end{tikzpicture}

\omcaption{The two markets lines banks report and the desks within them, with the quantitative jobs each desk tends to
employ (chapter~18 describes the bank-quant roles). Banks draw the lines differently; this is the common
shape.}\label{fig:in:bank-markets-divisions:tree}
\end{omfigure}

\section{American banks}

The five largest US banks report their markets lines in their annual reports on Form 10-K, three years at a time. In
2025 their equities shares ranged from 26\% to 64\%
(\cref{fig:in:bank-markets-divisions:mix}), and the share rose at each of the four that report 2023 figures, by 4.5 to
7.7 percentage points, as equities revenue grew faster than fixed income. The two banks whose businesses are
investment banking and markets rather than consumer banking earn more than half of their markets revenue in
equities; the largest consumer-and-corporate bank earns about three dollars in eight.

\begin{dated}{2025-12}{Nine banks' markets revenue by product line}\label{dat:in:bank-markets-divisions:mix}
\footnotesize
\begin{tabular}{@{}lllrrr@{}}
\toprule
bank & lines as reported & currency & \omterm{def:in:bank-markets-divisions:ficc}{FICC} & equities & equities share\\
\midrule
Morgan Stanley & Fixed Income; Equity & \$m & 8\,716 & 15\,631 & 64.2\%\\
Societe Generale & Fixed Income and Currencies; Equities & \euro m & 2\,346 & 3\,634 & 60.8\%\\
Goldman Sachs & \omterm{def:in:bank-markets-divisions:ficc}{FICC}; Equities & \$m & 14\,522 & 16\,535 & 53.2\%\\
BNP Paribas & \omterm{def:in:bank-markets-divisions:ficc}{FICC}; Equity \& Prime Services & \euro m & 5\,522 & 4\,046 & 42.3\%\\
Bank of America & \omterm{def:in:bank-markets-divisions:ficc}{FICC}; Equities (sales and trading) & \$m & 12\,267 & 8\,604 & 41.2\%\\
Barclays & \omterm{def:in:bank-markets-divisions:ficc}{FICC}; Equities & \pounds m & 5\,429 & 3\,225 & 37.3\%\\
JPMorgan Chase & Fixed Income Markets; Equity Markets & \$m & 22\,532 & 13\,250 & 37.0\%\\
Citigroup & Fixed Income Markets; Equity Markets & \$m & 16\,226 & 5\,744 & 26.1\%\\
Deutsche Bank & Fixed Income \& Currencies (incl.\ financing) & \euro m & 9\,600 & -- & 0\%\\
\bottomrule
\end{tabular}

\smallskip\noindent Full-year 2025 figures from the banks' Forms 10-K and results releases, ordered by equities share.
The US banks report 2023 figures too: equities shares were 31.4\% (JPMorgan Chase), 48.7\% (Goldman Sachs), 56.5\%
(Morgan Stanley) and 21.6\% (Citigroup).
\end{dated}

\section{The French banks and equity derivatives}

Two French banks sit on either side of the middle: one earns 61\% of its markets revenue in equities, the other 42\%.
The first describes its wholesale bank as offering ``distinctive global leadership in equity derivatives, structured
finance and ESG''. Equity derivatives in this sense are largely structured products (Book~5, chapter~19): notes whose
payoff depends on share indices and single shares, sold through networks to savers and institutions, and hedged by the
bank's trading desks (Book~9, chapter~28). A business of this kind employs many quants per dollar of revenue:
structurers design the products, library quants build the pricers, desk strategists hedge the exotic books, validators
check the models (chapter~18).

\begin{remark}[Why specialisation persists]\label{rem:in:bank-markets-divisions:why}
A bank's product strength is sticky. Pricing libraries, trained staff, distribution relationships and the risk of the
existing book take years to build and cannot be moved quickly; a bank strong in one product tends to stay so, and one
that exits a product rarely returns. The published numbers show the mix, not its causes; the causes the banks give are
in their own descriptions, which are statements of strategy, not evidence.
\end{remark}

\section{Other European banks}

A British bank earns 37\% of its markets revenue in equities, and a German bank none: having announced in July 2019
that it would ``exit its Equities Sales \& Trading business, while retaining a focused equity capital markets
operation'', it reported fixed income and currencies revenue of \euro9.6 billion in 2025, including financing
activities. An exit of that kind closes a line of quant jobs at one employer at once, and the people move to the
banks and firms that stay. Converted at the European Central Bank's average rates for 2025, the nine banks' markets
revenue totalled \$173.9 billion, of which \$72.7 billion was equities; the three most equity-heavy banks earned 49.9\%
of that equities revenue.

\begin{omfigure}
\begin{tikzpicture}
\begin{axis}[width=11cm,height=6.2cm,xbar stacked,bar width=9pt,xmin=0,xmax=40,xlabel={\small 2025 markets revenue (\$ billion, ECB average rates)},
  ytick={0,1,2,3,4,5,6,7,8},yticklabels from table={figdata/industry/07-bank-markets-divisions/mix2025.csv}{bank},y dir=reverse,
  tick label style={font=\footnotesize},grid=major,grid style={gray!20},table/col sep=comma,
  legend style={legend cell align=left,font=\scriptsize,at={(0.5,-0.2)},anchor=north,/tikz/every even column/.append style={column sep=8pt}},legend columns=2]
\addplot[fill=omThm!55,draw=omThm] table[y=k,x=equities_bn]{figdata/industry/07-bank-markets-divisions/mix2025.csv};
\addplot[fill=omDef!55,draw=omDef] table[y=k,x=ficc_bn]{figdata/industry/07-bank-markets-divisions/mix2025.csv};
\legend{equities,fixed income (FICC)}
\end{axis}
\end{tikzpicture}

\omcaption{Nine banks' 2025 markets revenue, equities and fixed income, converted to US dollars at the ECB's annual
average rates, ordered by equities share (not by size). Each bank's lines are its own and are not strictly comparable
in amount. Data: Forms 10-K and results releases for 2025, through \texttt{in\_banks.usd}.}\label{fig:in:bank-markets-divisions:mix}
\end{omfigure}

\section{League tables: who compiles them and what they measure}

\begin{definition}[Revenue pool, league table]\label{def:in:bank-markets-divisions:league}
A \emph{revenue pool}\index{revenue pool} is the total revenue that all dealers together earn from a product or a group
of clients over a period, as estimated by an analyst. A \emph{league table}\index{league table} ranks dealers by their
share of a revenue pool, of volumes or of clients' business, as estimated and published by an analytics firm, an
industry publication or a data vendor, for a stated period.
\end{definition}

Banks quote \omterm{def:in:bank-markets-divisions:league}{league tables} in their presentations, and candidates read them as a ranking of employers. They are neither
filings nor audited: they are estimates built by analytics firms from the banks' own submissions and from client
surveys. One such firm describes its service as helping banks ``measure revenue growth, wallet share, and retention
across sales, product and service'', with ``leading global and regional corporate investment banks and dealers'' as its
clients. A \omterm{def:in:bank-markets-divisions:league}{league table}'s rank depends on its definitions --- which products, which regions, which period --- and on
who submits data. This book cites no ranking. Where a table is useful, it states its publisher, its date and what it
measures, and a candidate should do the same.

\section{Quant headcounts: what is and is not public}

No bank in this chapter reports how many quants it employs. Annual reports give total headcount (chapter~1: 318\,512 at
one bank, a headcount of 47\,400 at another) and sometimes the number in the corporate and investment bank, but not the
number in pricing, risk or validation. The public evidence is indirect: the labour-condition filings of chapter~14 show
the bank-quant titles and their pay ranges by employer, and the remuneration disclosures of chapter~14 count the
material risk takers by business area. A candidate who wants to know the size of a team asks the team.

\section{Tutorial: where the equity derivatives are}\label{tut:in:bank-markets-divisions:tut}

\begin{tutorial}
\textbf{Goal.} Tabulate nine banks' markets revenue by product line from their filings, convert it to one currency and
order the banks by mix. \textbf{End state:} the dated box and \cref{fig:in:bank-markets-divisions:mix}.
\begin{tutsteps}
\item \textbf{The rows.} \texttt{data/industry/bank\_markets\_revenue.csv} holds one row per bank and year: the two
lines as the bank reports them, the currency, the ledger row that quotes the filing, and the bank's own line names.
\item \textbf{The rates.} \texttt{data/industry/ecb\_fx\_annual.csv} holds the ECB's annual average reference rates;
\texttt{firm.bankmix.load\_fx} turns them into US dollars per euro and per pound.
\item \textbf{The mix.} \texttt{firm.bankmix.mix\_table(rows, 2025)} orders the banks by equities share
(\cref{lst:in:bank-markets-divisions:mix}); \texttt{in\_banks.changes()} gives each bank's change in share between its
first and last reported years.
\omcode{code/firm/bankmix/firm_bankmix.py}{52}{70}{The equities share, the ordering by mix, and the share of equities
revenue earned by the most equity-heavy banks.}\label{lst:in:bank-markets-divisions:mix}
\item \textbf{Concentration.} \texttt{in\_banks.top3()} converts equities revenue to dollars and computes the share of
the three most equity-heavy banks.
\end{tutsteps}
\end{tutorial}

\textbf{What to change next.} Add a bank from its own results release, with its own line names (exercise~7); recompute
the dollar totals at 2024 rates to see how much of a year's change is currency.

\section{Build: the markets mix}\label{bld:in:bank-markets-divisions:bankmix}

\begin{build}
\bfield{Purpose} Put banks' markets revenue in two standard lines and one currency, so that chapters~12 and~18 can read a
bank's product strength beside its pay data and roles.
\bfield{Interface} \texttt{firm.bankmix}: \texttt{Row(bank, year, ficc, equities, currency, ledger, lines)};
\texttt{load(path)}; \texttt{load\_fx(path)}; \texttt{to\_usd(row, fx)}; \texttt{equities\_share(row)};
\texttt{mix\_table(rows, year)}; \texttt{mix\_change(rows, bank, y0, y1)}; \texttt{top\_k\_equities\_share(rows, fx, year, k)}.
\bfield{Rules} Every row keeps the bank's own line names and its source; amounts are converted at the year's average
rate, never at a spot rate; banks are ordered by mix, never by size alone or by any ranking.
\bfield{Acceptance tests} \texttt{code/firm/bankmix/tests/}: shares and order on constructed rows; the change in share in
percentage points; conversion through the euro cross rate; the top-$k$ share; a CSV round trip.
\bfield{Stretch} Quarterly rows and the seasonality of the mix; a regression of each bank's equities share on the VIX
index's yearly mean (Book~16's table).
\end{build}

\begin{omsources}
\item Forms 10-K for 2025 of JPMorgan Chase, Goldman Sachs, Morgan Stanley, Citigroup and Bank of America.
\item Results releases for 2025: BNP Paribas (appendices, 5 February 2026), Societe Generale (6 February 2026), Barclays
(2025 Results Announcement), Deutsche Bank (29 January 2026); Deutsche Bank media release of 7 July 2019.
\item European Central Bank, euro foreign exchange reference rates, annual averages.
\item Crisil Coalition Greenwich, description of its benchmarking services.
\end{omsources}

\section{Exercises}

\begin{exercise}[$\star$]\label{exo:in:bank-markets-divisions:1}
Compute Citigroup's equities share of markets revenue in 2023 and 2025, and its change in percentage points.
\end{exercise}

\begin{exercise}[$\star$]\label{exo:in:bank-markets-divisions:2}
Convert the British bank's 2025 \omterm{def:in:bank-markets-divisions:ficc}{FICC} income of \pounds5\,429 million to dollars at the ECB's average rates.
\end{exercise}

\begin{exercise}[$\star$]\label{exo:in:bank-markets-divisions:3}
Which desks of \cref{fig:in:bank-markets-divisions:tree} would you expect to employ the most model quants at a bank
whose equities share is 60\%, and why?
\end{exercise}

\begin{exercise}[$\star\star$]\label{exo:in:bank-markets-divisions:4}
Morgan Stanley's equities share rose from 56.5\% to 64.2\% between 2023 and 2025. Decompose the change: by how much did
equities and fixed-income revenue each grow?
\end{exercise}

\begin{exercise}[$\star\star$]\label{exo:in:bank-markets-divisions:5}
Why should the amounts of \cref{fig:in:bank-markets-divisions:mix} not be read as a ranking of the banks' markets
businesses?
\end{exercise}

\begin{exercise}[$\star\star$]\label{exo:in:bank-markets-divisions:6}
A \omterm{def:in:bank-markets-divisions:league}{league table} ranks a bank first in a product with a 15\% share of the \omterm{def:in:bank-markets-divisions:league}{revenue pool}. List four questions to ask before
using that rank to compare employers.
\end{exercise}

\begin{exercise}[$\star\star\star$]\label{exo:in:bank-markets-divisions:7}
\emph{Coding.} Add the German bank as a row with zero equities and recompute the three most equity-heavy banks' share of
equities revenue if the German bank had kept an equities business earning \euro3 billion. What changes, and why?
\end{exercise}

\begin{exercise}[$\star\star\star$]\label{exo:in:bank-markets-divisions:8}
\emph{Find the flaw.} ``The French bank earns 61\% in equities and the US bank 26\%: equities quants are paid more at the
French bank.''
\end{exercise}

\section{Problem: Where the Equity Derivatives Are}

\begin{problem}[{Weekend problem --- where the equity derivatives are}]\label{pb:in:bank-markets-divisions:1}
A derivatives quant with offers from a US bank and a European bank wants to know which has the stronger equities
business and how much weight to put on \omterm{def:in:bank-markets-divisions:league}{league tables}.

\medskip\noindent\textbf{Part I --- The lines.}
\begin{enumerate}
\item Define \omterm{def:in:bank-markets-divisions:ficc}{FICC} and the equities line, and name the desks of each.
\item Why are comparisons of shares within banks safer than comparisons of amounts across them?
\item Give the nine banks' 2025 equities shares in order.
\item Give the four US banks' 2023 equities shares and their changes to 2025.
\item What did the German bank announce in July 2019, and what does it report now?
\end{enumerate}

\medskip\noindent\textbf{Part II --- The conversion.}
\begin{enumerate}[resume]
\item Give the ECB's 2025 average rates and the implied dollars per pound.
\item Convert the French banks' 2025 equities revenue to dollars.
\item Give the nine banks' total markets revenue and equities revenue in dollars.
\item Give the share of equities revenue earned by the three most equity-heavy banks.
\item Why is an average rate used rather than the year-end rate?
\end{enumerate}

\medskip\noindent\textbf{Part III --- The jobs.}
\begin{enumerate}[resume]
\item Which quant roles does a bank strong in equity derivatives employ, and why many per dollar of revenue?
\item Why is a bank's product strength sticky?
\item What do banks publish about their quant headcounts?
\item Where can a candidate find indirect evidence of a bank's quant roles and pay?
\item What happens to a product's quant jobs when a bank exits it?
\end{enumerate}

\medskip\noindent\textbf{Part IV --- The verdict.}
\begin{enumerate}[resume]
\item State the \emph{named result}: the range of equities shares across the nine banks in 2025, and the three most
equity-heavy banks' share of equities revenue.
\item Define a \omterm{def:in:bank-markets-divisions:league}{revenue pool} and a \omterm{def:in:bank-markets-divisions:league}{league table}.
\item What must a cited \omterm{def:in:bank-markets-divisions:league}{league table} state?
\item Which of the quant's two offers is at the bank with the stronger equities mix, if the US bank is the one at 26\%
and the European one at 61\%?
\item In two sentences, what should the quant weigh beyond the mix?
\end{enumerate}
\end{problem}

\section{Interview questions}

\begin{interviewq}[$\star$ \iqroles{bank}]\label{iq:in:bank-markets-divisions:1}
What does \omterm{def:in:bank-markets-divisions:ficc}{FICC} stand for, and how does its revenue differ in kind from equities revenue?
\end{interviewq}

\begin{interviewq}[$\star$ \iqroles{bank, trader}]\label{iq:in:bank-markets-divisions:2}
Why might a bank's equities revenue grow faster than its fixed-income revenue in a year of rising share prices and
volatility?
\end{interviewq}

\begin{interviewq}[$\star\star$ \iqroles{bank}]\label{iq:in:bank-markets-divisions:3}
A bank sells structured notes to retail savers through its branches. Which desks and quants does that business need,
from product design to hedging?
\end{interviewq}

\begin{interviewq}[$\star\star$ \iqroles{bank, risk}]\label{iq:in:bank-markets-divisions:4}
Why would a bank exit equities sales and trading altogether rather than shrink it?
\end{interviewq}

\begin{interviewq}[$\star\star$ \iqroles{researcher}]\label{iq:in:bank-markets-divisions:5}
You have four years of two revenue lines for nine banks. How would you test whether the mix is more stable than the
level?
\end{interviewq}

\begin{interviewq}[$\star\star\star$ \iqroles{bank}]\label{iq:in:bank-markets-divisions:6}
Your equity-derivatives desk's revenue rose 30\% in a year. What part of that would you attribute to the market, what to
the franchise, and how would you tell?
\end{interviewq}
